\chapter{El colectivo macrocanónico}\label{ch:b3:grand-canonical}

Una gota de agua en aire húmedo ni crece ni se encoge; el oxígeno se une a
la hemoglobina en los pulmones y se suelta en los músculos; una pila empuja
electrones a través de un teléfono porque son ``más caros'' en un electrodo
que en el otro. En cada caso, la moneda que se equilibra no es la
temperatura, sino el \emph{\omterm{thm:b3:grand-canonical:distribution}{potencial químico}} $\mu$: el coste en \omterm{prop:b3:canonical-ensemble:machine}{energía
libre} de una partícula más. Este capítulo abre la segunda válvula del
reservorio --- partículas además de energía --- y construye el \omterm{thm:b3:grand-canonical:distribution}{colectivo
macrocanónico}, cuya recompensa es una eficacia asombrosa: la ocupación de un
único nivel cuántico, sea fermiónico o bosónico, se sigue en dos líneas, y
con ella las fórmulas maestras del capítulo siguiente. Sus aplicaciones van
de las isotermas de adsorción de las máscaras antigás y los catalizadores a
la \omterm{prop:b3:grand-canonical:saha}{ecuación de Saha}, con la que una estudiante de doctorado de Harvard leyó
en 1925 las temperaturas de las estrellas y descubrió, contra toda
expectativa, que el universo está hecho de hidrógeno.

\section{Las partículas entran en la negociación}

\begin{theorem}[Distribución macrocanónica]\label{thm:b3:grand-canonical:distribution}
\index{colectivo macrocanónico}\index{potencial químico}
Un sistema que intercambia energía \emph{y partículas} con un reservorio a
temperatura $T$ y \omterm{thm:b3:grand-canonical:distribution}{potencial químico} $\mu$ ocupa su estado $s$
(de energía $E_s$ y número de partículas $N_s$) con probabilidad
\[
\mathcal P_s = \frac{\eu^{-\beta(E_s - \mu N_s)}}{\Xi} , \qquad
\Xi = \sum_s\eu^{-\beta(E_s - \mu N_s)} ,
\]
la \emph{\omterm{thm:b3:canonical-ensemble:boltzmann}{función de partición} macrocanónica}. De $\Xi$: $\langle N\rangle
= k_{\text{B}}T\,\partial\ln\Xi/\partial\mu$, y el \emph{potencial
macrocanónico} $\Phi = -k_{\text{B}}T\ln\Xi$ cumple $\Phi = -PV$.
\end{theorem}

\begin{proof}[Demostración parcial]
Como en el \cref{thm:b3:canonical-ensemble:boltzmann}, desarróllese la
entropía del reservorio en las dos retiradas:
$S_{\text{res}}(E - E_s, N - N_s) \approx S_{\text{res}} - E_s/T +
\mu N_s/T$, usando $\partial S/\partial N = -\mu/T$
(\cref{thm:b3:microcanonical-ensemble:temperature}). Los promedios
se siguen por derivación; $\Phi = -PV$ se admite (se sigue de la
extensividad).
\end{proof}

\begin{definition}[Qué significa $\mu$]\label{def:b3:grand-canonical:mu}
\index{potencial químico!significado}
$\mu$ es el coste, en \omterm{prop:b3:canonical-ensemble:machine}{energía libre}, de añadir una partícula a $T$ y $V$
fijos. Las partículas fluyen espontáneamente de $\mu$ alto a $\mu$
bajo --- la difusión, la ósmosis, la evaporación, las reacciones químicas y
las corrientes eléctricas son todas gradientes de $\mu$ que se relajan. En
el equilibrio entre dos fases o dos lugares cualesquiera, $\mu$ es igual a
ambos lados --- la tercera igualdad (junto con $T$ y $P$) que gobierna la
coexistencia. Para un gas ideal clásico,
\[
\mu = k_{\text{B}}T\,\ln\!\big(n\lambda_T^3\big) < 0
\quad\text{(gas diluido)} :
\]
añadir una partícula a un gas diluido \emph{baja} la \omterm{prop:b3:canonical-ensemble:machine}{energía libre},
porque la ganancia de entropía supera cualquier coste energético --- el
signo que impulsa la evaporación hacia el aire seco. Y una pila es una
máquina de $\mu$: su tensión es la caída de \omterm{thm:b3:grand-canonical:distribution}{potencial químico} por electrón
entre sus electrodos, $U = \Delta\mu/e$ --- los voltios son el tipo de
cambio de la química del electrón, y por eso las pilas entregan voltajes de
una sola cifra.
\end{definition}

\section{Un nivel cada vez: el truco maestro}

\begin{theorem}[Ocupación de un solo nivel cuántico]\label{thm:b3:grand-canonical:occupation}
\index{distribución de Fermi--Dirac}\index{distribución de Bose--Einstein}
Tómese como ``el sistema'' un único nivel de una partícula, de energía
$\epsilon$, en contacto macrocanónico con el gas de todos los demás.
\emph{Fermiones} (ocupación $0$ o $1$):
\[
\Xi = 1 + \eu^{-\beta(\epsilon - \mu)}
\;\Longrightarrow\;
\bar n_{\text{FD}} = \frac{1}{\eu^{\beta(\epsilon - \mu)} + 1} .
\]
\emph{Bosones} (ocupación $0, 1, 2, \dots$):
\[
\Xi = \frac{1}{1 - \eu^{-\beta(\epsilon - \mu)}}
\;\Longrightarrow\;
\bar n_{\text{BE}} = \frac{1}{\eu^{\beta(\epsilon - \mu)} - 1} ,
\]
con la exigencia $\mu < \epsilon_{\min}$. Ambas se reducen al factor de
Boltzmann $\eu^{-\beta(\epsilon - \mu)}$ cuando las ocupaciones son escasas.
Estas dos fórmulas, obtenidas en dos líneas cada una, son todo el
material de partida de la estadística cuántica: el capítulo siguiente es su
cosecha.
\end{theorem}

\begin{proof}
Fermiones: dos términos. Bosones: una serie geométrica, convergente solo
para $\mu < \epsilon$. En ambos casos, $\bar n = k_{\text{B}}T\,
\partial\ln\Xi/\partial\mu$ da las formas enunciadas.
\end{proof}

\begin{omfigure}
\begin{tikzpicture}
  \begin{axis}[omaxis, axis x line=bottom, axis y line=left,
    width=11.2cm, height=5.6cm, xmin=-6, xmax=6, ymin=0, ymax=2.3,
    xlabel={$(\epsilon - \mu)/k_{\text{B}}T$}, ylabel={$\bar n$},
    xlabel style={at={(axis description cs:0.5,-0.1)}, anchor=north},
    xtick={-4,-2,0,2,4}, ytick={0.5,1,2}, clip=false,
    legend style={font=\scriptsize, at={(0.97,0.93)}, anchor=north east,
      draw=none, fill=none}]
    \addplot[omDef, very thick, domain=-6:6, samples=200]
      {1/(exp(x)+1)};
    \addlegendentry{Fermi--Dirac}
    \addplot[omThm, very thick, domain=0.45:6, samples=200]
      {1/(exp(x)-1)};
    \addlegendentry{Bose--Einstein}
    \addplot[omExo, dashed, thick, domain=-0.8:6, samples=100]
      {exp(-x)};
    \addlegendentry{Boltzmann}
    \draw[gray, dashed] (axis cs:0,0) -- (axis cs:0,2.3);
  \end{axis}
\end{tikzpicture}

{\small Las tres leyes de ocupación frente a $(\epsilon -
\mu)/k_{\text{B}}T$: los fermiones saturan en uno por estado (un
escalón suavizado en $\epsilon = \mu$), los bosones se amontonan sin límite
cuando $\epsilon \to \mu$, y ambas se unen a la cola diluida de Boltzmann
por la derecha --- el régimen clásico es aquel en el que un estado apenas
llega a visitarse.}
\end{omfigure}

\section{Adsorción: sitios en alquiler}

\begin{proposition}[La isoterma de Langmuir]\label{prop:b3:grand-canonical:langmuir}
\index{isoterma de Langmuir}\index{adsorción}
Una superficie ofrece sitios independientes, cada uno vacío o con una
molécula de gas de energía de enlace $-\epsilon_0$; el gas de encima, a
presión $P$, es el reservorio. Cada sitio es un sistema macrocanónico de dos
estados, y su ocupación es
\[
\theta = \frac{1}{\eu^{\beta(-\epsilon_0 - \mu)} + 1}
 = \frac{P}{P + P_0(T)} , \qquad
P_0(T) = \frac{k_{\text{B}}T}{\lambda_T^3}\,
\eu^{-\epsilon_0/k_{\text{B}}T} ,
\]
usando el $\mu$ del gas ideal. La cobertura crece linealmente a baja
presión y satura en una monocapa --- la \emph{\omterm{prop:b3:grand-canonical:langmuir}{isoterma de Langmuir}}, la
curva de trabajo de los catalizadores de automóvil, el carbón activo, las
máscaras antigás y la fijación del oxígeno por la sangre
(\cref{exo:b3:grand-canonical:6}). La presión de cruce $P_0$
cae \emph{exponencialmente} con la energía de enlace y sube de forma abrupta
con la temperatura: calentar una superficie sacude a sus huéspedes hasta
soltarlos --- por eso hornear el material de vidrio al vacío es la única
manera de limpiarlo de verdad.
\end{proposition}

\begin{proof}
La ocupación del sitio es la del
\cref{thm:b3:grand-canonical:occupation} (en forma fermiónica: un sitio
alberga como mucho uno) con $\epsilon = -\epsilon_0$; sustitúyase
$\eu^{\beta\mu} = n\lambda_T^3 = P\lambda_T^3/k_{\text{B}}T$.
\end{proof}

\begin{omfigure}
\begin{tikzpicture}
  \begin{axis}[omaxis, axis x line=bottom, axis y line=left,
    width=10.8cm, height=5.4cm, xmin=0, xmax=5, ymin=0, ymax=1.12,
    xlabel={$P/P_0$ a la temperatura menor}, ylabel={cobertura $\theta$},
    xlabel style={at={(axis description cs:0.5,-0.1)}, anchor=north},
    xtick={1,2,3,4}, ytick={0.5,1}, clip=false,
    legend style={font=\scriptsize, at={(0.97,0.42)}, anchor=north east,
      draw=none, fill=none}]
    \addplot[omDef, very thick, domain=0:5, samples=120] {x/(x+1)};
    \addlegendentry{superficie fría}
    \addplot[omThm, very thick, domain=0:5, samples=120] {x/(x+6)};
    \addlegendentry{superficie caliente ($P_0$ mayor)}
    \draw[gray, dashed] (axis cs:0,1) -- (axis cs:5,1);
    \node[font=\scriptsize, anchor=south east, gray] at (axis cs:4.9,1.0)
      {una monocapa};
  \end{axis}
\end{tikzpicture}

{\small \omterm{prop:b3:grand-canonical:langmuir}{Isotermas de Langmuir}: la cobertura frente a la presión, saturando
en una monocapa. Caliéntese la superficie y la misma cobertura exigirá
mucha más presión --- la adsorción es una negociación de $\mu$ que la
superficie caliente no deja de perder.}
\end{omfigure}

\section{Equilibrios de ionización: la ecuación de Saha}

\begin{proposition}[La ecuación de Saha]\label{prop:b3:grand-canonical:saha}
\index{ecuación de Saha}
En un gas caliente, la reacción $\text{H} \rightleftharpoons \text{p} +
\text{e}$ se equilibra cuando $\mu_{\text{H}} = \mu_{\text{p}} +
\mu_{\text{e}}$. Escribir cada $\mu$ con la forma del gas ideal (con el
átomo acreditado con su energía de enlace $-E_{\text{I}}$) da
\[
\frac{n_{\text{e}}\,n_{\text{p}}}{n_{\text{H}}}
 = \frac{1}{\lambda_{T,\text{e}}^3}\,
\eu^{-E_{\text{I}}/k_{\text{B}}T}
\qquad\Big(\lambda_{T,\text{e}} =
h/\sqrt{2\pi m_{\text{e}}k_{\text{B}}T}\Big)
\]
(salvo factores de degeneración de espín del orden de la unidad). De nuevo
dos monedas que compiten: el factor de Boltzmann grava la ionización con
$E_{\text{I}} = \qty{13.6}{eV}$, mientras que el prefactor --- la
\emph{entropía del electrón liberado}, con sus $\lambda_T^{-3}$ estados
accesibles --- la subvenciona. La subvención es tan grande a las densidades
estelares que el hidrógeno se ioniza apreciablemente ya cerca de los
\qty{10000}{K}, donde $k_{\text{B}}T$ vale solo
\qty{0.9}{eV}: los equilibrios se disputan en \omterm{prop:b3:canonical-ensemble:machine}{energía libre}, no en energía
--- y los espectros de las estrellas son el libro de cuentas
(\cref{pb:b3:grand-canonical:1}).
\end{proposition}

\begin{proof}[Demostración parcial]
Hágase $\mu_{\text{H}} = \mu_{\text{p}} + \mu_{\text{e}}$ con $\mu_i
= k_{\text{B}}T\ln(n_i\lambda_{T,i}^3) + \varepsilon_i$
($\varepsilon_{\text{H}} = -E_{\text{I}}$, y los demás nulos);
exponénciese. Las longitudes de onda térmicas del protón y del átomo casi se
cancelan (masas iguales); los factores de degeneración se admiten.
\end{proof}

\begin{omfigure}
\begin{tikzpicture}
  \begin{axis}[omaxis, axis x line=bottom, axis y line=left,
    width=10.8cm, height=5.4cm, xmin=3000, xmax=25000, ymin=0, ymax=1.1,
    xlabel={$T$ (K)}, ylabel={fracción ionizada de hidrógeno},
    xlabel style={at={(axis description cs:0.5,-0.1)}, anchor=north},
    xtick={5000,10000,15000,20000}, ytick={0.5,1}, clip=false]
    \addplot[omDef, very thick, domain=3000:25000, samples=300]
      {1/(1+exp(-(x-9600)/1100))};
    \draw[gray, dashed] (axis cs:9600,0) -- (axis cs:9600,0.5);
    \node[font=\scriptsize, anchor=west] at (axis cs:10000,0.28)
      {medio ionizado cerca de \qty{10000}{K}};
  \end{axis}
\end{tikzpicture}

{\small La ionización de Saha del hidrógeno a densidad fotosférica
(esquemática): una transición brusca cerca de los \qty{10000}{K}, muy por
debajo de los ingenuos $E_{\text{I}}/k_{\text{B}} \approx \qty{160000}{K}$
--- la entropía del electrón liberado paga la mayor parte de la factura.}
\end{omfigure}

\begin{method}[El oficio macrocanónico]\label{met:b3:grand-canonical:method}
(1) Todo lo que intercambie partículas con un socio grande: dese al
socio un $\mu$ y pésense los estados con $\eu^{-\beta(E - \mu N)}$.
(2) Los sitios o niveles independientes se factorizan: resuélvase \emph{uno}
y multiplíquese. (3) Para un reservorio de gas clásico, $\eu^{\beta\mu} =
n\lambda_T^3$ convierte $\mu$ en presión medible. (4) Reacciones
y cambios de fase: equilíbrense los $\mu$, uno por especie, con las energías
de enlace incluidas --- espérese que los prefactores entrópicos desplacen
los equilibrios muy lejos de las conjeturas ingenuas de Boltzmann. (5)
Electrones en sólidos: su $\mu$ es el nivel de Fermi, y las diferencias de
$\mu$ son tensiones --- el puente hacia el
\cref{ch:b3:electrons-in-solids}.
\end{method}

\section{Ejercicios}

\begin{exercise}[$\star$]\label{exo:b3:grand-canonical:1}
(a) Se conectan dos cajas del mismo gas a la misma $T$ pero con densidades
distintas: ¿en qué sentido fluyen las partículas, en lenguaje de
$\mu$? (b) Muéstrese, a partir del $\mu$ del gas ideal, que el flujo cesa
con densidades iguales. (c) Hay agua bajo aire húmedo: enúnciese la
condición de equilibrio en términos de $\mu$. (d) ¿Por qué se seca la ropa
incluso por debajo de $100\,^\circ$C cuando el aire está seco (compárense los
$\mu$, no las temperaturas)?
\end{exercise}

\begin{exercise}[$\star$]\label{exo:b3:grand-canonical:2}
El \omterm{thm:b3:grand-canonical:distribution}{potencial químico} del aire. (a) Calcúlese $\lambda_T$ para el N$_2$
a \qty{300}{K}. (b) Calcúlese $n\lambda_T^3$ a \qty{1}{bar} y
el $\mu$ resultante en eV. (c) ¿Por qué no es paradójico $\mu < 0$
(qué término de $F = E - TS$ domina al añadir una molécula)? (d)
¿A qué densidad alcanzaría $\mu$ el valor cero, y qué anuncia
$n\lambda_T^3 \sim 1$
(\cref{exo:b3:grand-canonical:11})?
\end{exercise}

\begin{exercise}[$\star$]\label{exo:b3:grand-canonical:3}
Dedúzcase la ocupación de Fermi--Dirac a partir de su $\Xi$ de dos términos,
paso a paso; compruébense los límites $\epsilon \ll \mu$ y $\epsilon \gg
\mu$, y el valor en $\epsilon = \mu$. ¿Dónde ha aparecido ya esta función
en este libro
(\cref{exo:b3:microcanonical-ensemble:8})?
\end{exercise}

\begin{exercise}[$\star$]\label{exo:b3:grand-canonical:4}
Los voltios son química. (a) Una pila de litio da \qty{3.0}{V}:
exprésese la diferencia de \omterm{thm:b3:grand-canonical:distribution}{potencial químico} por electrón en eV y
por mol en kJ. (b) ¿Por qué tienen todas las pilas cotidianas voltajes de
una sola cifra (qué fija la escala de las diferencias químicas de $\mu$)? (c) Una
batería de plomo y ácido de coche apila seis celdas: ¿por qué apilar en vez
de buscar una reacción de \qty{12}{eV}? (d) En lenguaje de $\mu$, ¿qué hace
una pila cuando está ``gastada''?
\end{exercise}

\begin{exercise}[$\star\star$]\label{exo:b3:grand-canonical:5}
Bosones en un nivel. (a) Súmese la serie geométrica de $\Xi$ y
dedúzcase $\bar n_{\text{BE}}$. (b) ¿Por qué debe quedarse $\mu$ por debajo
del nivel más bajo --- qué diverge si no, y qué anticipa esa
divergencia (\cref{ch:b3:quantum-statistics})? (c)
Calcúlese $\bar n$ para $(\epsilon - \mu)/k_{\text{B}}T = 0.1$, $1$ y
$3$, y compárese con Boltzmann. (d) Los fotones tienen $\mu = 0$: dese
la razón física (su número no se conserva --- las paredes los absorben y los
emiten sin trabas).
\end{exercise}

\begin{exercise}[$\star\star$]\label{exo:b3:grand-canonical:6}
Langmuir en acción. (a) A partir de la isoterma, ¿qué presión logra
$\theta = 0.5$? ¿Y $0.9$? (b) La mioglobina une O$_2$ en un solo sitio:
la saturación sigue a Langmuir en la presión parcial de oxígeno ---
¿por qué se observa una curva \emph{hiperbólica} (y no sigmoidal) para la
mioglobina y una curva cooperativa en S para la hemoglobina, de cuatro
sitios? (c) El monóxido de carbono se une a los mismos sitios con un
$\epsilon_0$ mayor en $\approx\qty{0.13}{eV}$: a presiones parciales iguales
y \qty{310}{K}, calcúlese la razón de ocupaciones CO:O$_2$ y explíquese el
envenenamiento por CO. (d) ¿Por qué funciona la oxigenoterapia hiperbárica,
en lenguaje de isotermas?
\end{exercise}

\begin{exercise}[$\star\star$]\label{exo:b3:grand-canonical:7}
Fluctuaciones del número. (a) Muéstrese que $\operatorname{Var}N =
k_{\text{B}}T\,\partial\langle N\rangle/\partial\mu$. (b) Para el
gas ideal clásico, dedúzcase $\operatorname{Var}N = \langle
N\rangle$: Poisson. (c) Reconcíliese con las fluctuaciones de densidad del
\cref{exo:b3:microcanonical-ensemble:12}. (d) Para un solo nivel
fermiónico, calcúlese $\operatorname{Var}n = \bar n(1 - \bar
n)$: ¿dónde están más silenciosos los números de fermiones, y por qué?
\end{exercise}

\begin{exercise}[$\star\star$]\label{exo:b3:grand-canonical:8}
La ósmosis a partir de $\mu$. Una membrana deja pasar el agua pero no el
soluto (de concentración $c$ por volumen, diluido). (a) Igualando el $\mu$
del agua a los dos lados de la membrana, muéstrese que el lado de agua pura
debe estar sobrepresurizado en $\Pi = ck_{\text{B}}T$ (van 't Hoff) ---
acéptese que el soluto disuelto baja el $\mu$ del agua en
$-k_{\text{B}}T\,c/n_{\text{w}}$ por molécula de agua. (b) Calcúlese
$\Pi$ para una solución salina fisiológica ($c \approx \qty{3e26}{m^{-3}}$
de partículas totales). (c) Glóbulos rojos en agua pura: ¿en qué sentido
fluye el agua y con qué resultado? (d) ¿Por qué no necesitan los árboles ni
bombas ni milagros para subir la savia decenas de metros (gradientes
osmóticos y capilares de $\mu$)?
\end{exercise}

\begin{exercise}[$\star\star$]\label{exo:b3:grand-canonical:9}
Equilibrio en un campo. En un campo gravitatorio, el equilibrio de una
columna de gas exige que el \omterm{thm:b3:grand-canonical:distribution}{potencial químico} \emph{total} $\mu(h) =
k_{\text{B}}T\ln(n\lambda_T^3) + mgh$ sea uniforme. (a) Dedúzcase de aquí
la ley barométrica en una línea. (b) Generalícese a una centrifugadora de
velocidad angular $\omega$ (con potencial
$-\tfrac12 m\omega^2r^2$) y obténgase el perfil radial. (c)
Las centrifugadoras de enriquecimiento de uranio giran a $\omega r \sim
\qty{600}{m/s}$: calcúlese la razón de equilibrio del enriquecimiento de
$^{238}$UF$_6$ frente a $^{235}$UF$_6$ por etapa, $\exp[\Delta m\,\omega^2
r^2/2k_{\text{B}}T]$ con $\Delta m = \qty{5e-27}{kg}$ y $T =
\qty{320}{K}$. (d) ¿Por qué se encadenan muchas etapas?
\end{exercise}

\begin{exercise}[$\star\star\star$]\label{exo:b3:grand-canonical:10}
Saha, a mano. (a) Realícese la demostración con el balance de $\mu$ de la
\cref{prop:b3:grand-canonical:saha}. (b) Defínase la fracción ionizada
$x$ a densidad total de hidrógeno $n$ y muéstrese que $x^2/(1 - x)
= \eu^{-E_{\text{I}}/k_{\text{B}}T}/n\lambda_{T,\text{e}}^3$. (c)
Fotosfera solar: $T = \qty{5800}{K}$, $n \approx
\qty{e23}{m^{-3}}$: muéstrese que la fracción ionizada de hidrógeno es solo
de $\sim\num{e-4}$ --- la superficie del Sol es neutra. (d) A $T =
\qty{10000}{K}$ (con la misma $n$): recalcúlese y coméntese lo
abrupto del cambio.
\end{exercise}

\begin{exercise}[$\star\star\star$]\label{exo:b3:grand-canonical:11}
Cuando los gases se vuelven cuánticos. La degeneración se instala en el
punto $n\lambda_T^3 \sim
1$. (a) Para el aire a \qty{300}{K}: ¿cuánto por debajo
del umbral (\cref{exo:b3:grand-canonical:2})? (b) Para los electrones de
conducción del cobre ($n = \qty{8.5e28}{m^{-3}}$): hállese la temperatura
por debajo de la cual $n\lambda_T^3 > 1$ --- y conclúyase que los electrones
de un metal están degenerados a \emph{cualquier} temperatura de laboratorio.
(c) Para átomos de rubidio a $n = \qty{e20}{m^{-3}}$: la temperatura
umbral cuántica (la escala del condensado de Bose--Einstein del
\cref{ch:b3:quantum-statistics}). (d) Ordénense los tres y enúnciese la
regla: las partículas ligeras y las densidades altas se vuelven cuánticas
primero.
\end{exercise}

\begin{exercise}[$\star\star\star$]\label{exo:b3:grand-canonical:12}
El universo pierde sus positrones. En el universo primitivo caliente,
$\eu^+ + \eu^- \rightleftharpoons 2\gamma$ mantenía los pares en
equilibrio con $\mu_{e^+} + \mu_{e^-} = 0$ (los fotones: $\mu = 0$).
(a) Para un plasma simétrico esto da $\mu_{e^\pm} \approx 0$:
muéstrese que la densidad de pares es entonces $n_\pm \approx
\lambda_T^{-3}\eu^{-m_{\text{e}}c^2/k_{\text{B}}T}$ una vez que
$k_{\text{B}}T \ll m_{\text{e}}c^2$. (b) Calcúlense la temperatura
a la que $m_{\text{e}}c^2 = k_{\text{B}}T$ y la escala de tiempo cósmico
que le corresponde (acéptese $t \sim \qty{1}{s}$ a
$\qty{e10}{K}$). (c) Muéstrese que un enfriamiento adicional modesto (un
factor de unas pocas unidades) desploma la densidad de pares en muchos
órdenes de magnitud: los positrones se aniquilan y dejan solo el diminuto
exceso primordial de electrones. (d) Los fotones liberados recalientan el
baño de radiación: ¿qué radiación reliquia lleva el recibo
(\cref{ch:b3:astrophysics})?
\end{exercise}

\section{Problema: leer las temperaturas de las estrellas}

\begin{problem}[{Problema de fin de semana --- Saha, Payne y de qué está
hecho el universo}]\label{pb:b3:grand-canonical:1}
Los espectros estelares se catalogaron a lo largo de la década de 1890 en
las clases O, B, A, F, G, K y M según la intensidad de sus líneas de
absorción --- con las líneas de Balmer del hidrógeno máximas en la clase A y
misteriosamente débiles tanto en las estrellas más calientes como en las más
frías. En 1925, Cecilia Payne aplicó la flamante \omterm{prop:b3:grand-canonical:saha}{ecuación de Saha} y demostró
que todo el zoo era una misma sustancia a distintas temperaturas --- y que
las estrellas son abrumadoramente hidrógeno, conclusión tan herética que
tuvo que suavizarla en su tesis. Este problema recorre su razonamiento.
Datos: la absorción de Balmer parte de $n = 2$, a $E_2 - E_1 =
\qty{10.2}{eV}$ sobre el estado fundamental; $E_{\text{I}} =
\qty{13.6}{eV}$; densidades electrónicas fotosféricas $n_{\text{e}}
\sim \qty{e20}{m^{-3}}$.

\medskip\noindent\textbf{Parte I --- Dos factores en competencia.}
\begin{enumerate}
  \item La absorción de Balmer necesita hidrógeno neutro \emph{en} el nivel $n =
        2$: escríbase el factor de Boltzmann para la población de $n = 2$
        (degeneraciones: $g_2/g_1 = 4$).
  \item Evalúese a $5000$, $10000$ y \qty{20000}{K}: ¿en qué sentido
        empuja este factor la intensidad de la línea al subir $T$?
  \item Ahora el factor de Saha: al subir $T$, ¿qué le ocurre a la
        fracción \emph{neutra} y, por tanto, a las líneas de Balmer?
  \item Explíquese cualitativamente por qué el producto --- excitado
        \emph{y} todavía neutro --- debe tener un máximo a alguna
        temperatura intermedia.
  \item Usando la \omterm{prop:b3:grand-canonical:saha}{ecuación de Saha} con $n_{\text{e}} =
        \qty{e20}{m^{-3}}$, estímese la temperatura de
        semiionización del hidrógeno.
  \item Júntelo todo en el máximo: ¿cerca de qué temperatura deberían ser
        más intensas las líneas de Balmer, y a qué clase espectral
        (las estrellas A, $\sim\qty{10000}{K}$) señala eso?
\end{enumerate}

\medskip\noindent\textbf{Parte II --- La secuencia espectral,
descifrada.}
\begin{enumerate}[resume]
  \item Las estrellas M (\qty{3000}{K}) muestran bandas moleculares
        intensas (TiO) e hidrógeno débil: explíquense ambas cosas con los
        dos factores.
  \item Las estrellas O (\qty{40000}{K}) muestran líneas de helio y de
        nuevo hidrógeno débil: explíquese (qué le ha pasado al hidrógeno y
        por qué le llega ahora el turno al helio, con
        $E_{\text{I}} = \qty{24.6}{eV}$).
  \item El Sol (G, \qty{5800}{K}) muestra líneas intensas de
        \emph{calcio ionizado} ($E_{\text{I,Ca}} = \qty{6.1}{eV}$)
        y, sin embargo, hidrógeno débil: muéstrese con la lógica de Saha
        por qué el calcio ya está ionizado mientras que el hidrógeno sigue
        neutro y sobre todo en $n = 1$.
  \item Los astrónomos anteriores a Payne leían la intensidad de línea
        como abundancia: ¿qué les hicieron concluir sobre la composición
        del Sol sus poderosas líneas de calcio y sus enclenques líneas de
        Balmer?
  \item La intuición de Payne: corregir la intensidad de cada línea con su
        ``factor de visibilidad'' de Boltzmann y Saha. Cuando lo hizo,
        ¿qué jerarquía de abundancias verdaderas emergió?
  \item Su resultado --- hidrógeno un millón de veces más abundante
        que el calcio, y estrellas hechas sobre todo de hidrógeno y helio
        --- ¿por qué se le opuso resistencia (qué sugería la química
        centrada en la Tierra) y cuándo se la reivindicó?
\end{enumerate}

\medskip\noindent\textbf{Parte III --- Números sobre el máximo.}
\begin{enumerate}[resume]
  \item Calcúlese el factor de Boltzmann de $n = 2$ exactamente a
        \qty{9600}{K}.
  \item Con la estimación de Saha del apartado 5, tómese la fracción
        neutra a $9600$, $15000$ y \qty{25000}{K} como
        aproximadamente $0.5$, $\num{e-3}$ y $\num{e-5}$: fórmese el producto
        (fracción excitada $\times$ fracción neutra) a esas tres
        temperaturas y localícese su máximo.
  \item ¿Por qué es \emph{abrupto} el máximo de Balmer por el lado
        caliente (qué exponencial gana) y \emph{gradual} por el lado
        frío?
  \item La fotosfera de una enana blanca tiene una $n_{\text{e}}$ mil
        veces mayor: ¿en qué sentido se desplaza la temperatura de
        semiionización, y por qué (Le Chatelier en lenguaje de
        $\mu$)?
  \item Dos estrellas de idéntica temperatura pero distinta presión
        muestran, así, intensidades de línea distintas: ¿qué propiedad
        estelar (gigante frente a enana --- la gravedad superficial)
        permite eso leer a los espectroscopistas?
  \item Los cartografiados modernos fijan las temperaturas estelares con
        un $\pm1\%$ a partir de razones de líneas: enúnciese la cadena
        postulado $\to$ Boltzmann $\to$ Saha $\to$ termómetro.
\end{enumerate}

\medskip\noindent\textbf{Parte IV --- Qué quedó zanjado.}
\begin{enumerate}[resume]
  \item La secuencia OBAFGKM había sido un caos alfabético
        reordenado empíricamente: tras Saha, ¿qué única variable física
        la ordena?
  \item Dígase por qué el \emph{mismo} elemento puede dominar un
        espectro a una temperatura y desvanecerse a otra sin ningún
        cambio de abundancia --- la lección central que Payne enseñó a la
        astronomía.
  \item El hidrógeno, con tres cuartas partes de toda la masa bariónica:
        nómbrense otros dos pilares de este libro que descansan en esa
        abundancia (la fusión estelar, \cref{ch:b3:astrophysics}; el cielo
        de 21 cm, \cref{ch:b3:spin-two-level}).
  \item Otto Struve llamó a la tesis de Payne ``la tesis doctoral más
        brillante jamás escrita en astronomía'': a partir de este
        problema, justifíquese el elogio en una frase.
  \item La misma física de Saha, aplicada a la inversa a escala cósmica,
        predice la temperatura a la que se combinó por primera vez todo el
        hidrógeno del universo ($\sim\qty{3000}{K}$, a la densidad cósmica
        mucho menor): nómbrense el suceso y su reliquia
        (\cref{ch:b3:astrophysics}).
  \item Los equilibrios de ionización gobiernan también las lámparas
        fluorescentes y el procesado por plasma: dígase en una frase por qué
        puede existir un estado ionizado ``de tipo \qty{10000}{K}'' en un
        tubo cuyo vidrio permanece frío (¿qué dos subsistemas no llegan a
        equilibrarse?).
  \item Resúmase el resultado con nombre propio: una sola condición de
        equilibrio, $\mu_{\text{H}} = \mu_{\text{p}} + \mu_{\text{e}}$,
        convierte el alfabeto espectral en un termómetro que va de los
        $3000$ a los \qty{40000}{K}, sitúa el máximo de Balmer en las
        estrellas A y revela un universo hecho de hidrógeno.
\end{enumerate}
\end{problem}
